\documentclass[11pt]{article}
\usepackage[margin=0.86in]{geometry}
\usepackage[T1]{fontenc}
\usepackage[utf8]{inputenc}
\usepackage{amsmath,amssymb,amsthm}
\usepackage{graphicx}
\usepackage{xcolor}
\usepackage{microtype}
\usepackage[colorlinks=true,urlcolor=blue!55!black,citecolor=blue!55!black]{hyperref}

\newtheorem{theorem}{Theorem}
\newtheorem{lemma}[theorem]{Lemma}
\newcommand{\Torus}{\mathbb T}
\newcommand{\N}{\mathbb N}
\newcommand{\Ucal}{\mathcal U}

\title{Exact circle shape optimization for geodesic observation\\and for resonant wave times}
\author{Candidate partial answer to the shape-optimization problems of arXiv:1607.01535}
\date{11 August 2026}

\begin{document}
\maketitle

\begin{center}
\fbox{\parbox{0.92\textwidth}{\textbf{Status: candidate partial result, likely valid.}
On the circle, the finite-time geometric functional $g_2^T$ has supremum
$L$ for every observation time, with a complete attainment classification by
rational rotations.  At the resonant times $T\in2\pi\N$, the deterministic
wave constant satisfies $C_T=(T/2)g_1$ exactly; its supremum is $TL/2$ and it
is attained precisely at volume fraction $L=1/2$ by antipodal selectors.}}
\end{center}

\section*{Source problem}

Humbert, Privat, and Tr\'elat ask on PDF page 22 of
\emph{Observability properties of the homogeneous wave equation on a closed
manifold} \cite{source}:
\begin{quote}
A challenging problem is to maximize $\omega\mapsto C_T(\omega)$ over all
measurable subsets of prescribed volume fraction $L$.  Maximizing the
geometric functional $\omega\mapsto g_2(\omega)$ is also an interesting open
problem and would be a step toward the deterministic problem.
\end{quote}
The source passage is reproduced below.

\begin{center}
\includegraphics[width=0.96\textwidth]{figures/open_problem_crop.png}
\end{center}

We solve the complete circle subcase for the long-time functional $g_2$, and
more strongly for its finite-time version $g_2^T$.  We also solve the
deterministic problem at every common spectral period.  The arbitrary-time
deterministic problem and higher-dimensional manifolds remain open.

\section*{Definitions}

Write $\Torus=\mathbb R/(2\pi\mathbb Z)$, with arc-length measure, and
\[
 \Ucal_L=\{\omega\subset\Torus:\omega\text{ measurable and }
 |\omega|=2\pi L\},\qquad 0<L<1.
\]
Every unit-speed geodesic is $t\mapsto x\pm t$.  Hence the source's geometric
quantity is
\begin{equation}\label{eq:g2def}
 g_2^T(\omega)=\inf_{x\in\Torus}\frac1T
 \int_0^T\chi_\omega(x+t)\,dt.
\end{equation}
The sign of the ray is immaterial after changing $x$.

For the wave equation on the zero-mean sector used in the source's circle
example, let $C_T(\omega)$ be its deterministic observability constant and let
\[
 g_1(\omega)=\inf_{j\geq1}\ \inf_{0\neq\phi\in
 \operatorname{span}\{\cos(jx),\sin(jx)\}}
 \frac{\int_\omega|\phi|^2}{\int_\Torus|\phi|^2}.
\]

\section*{Proof intuition}

For $g_2^T$, remove the whole turns made by a geodesic.  What remains is the
least mass of $\omega$ in a sliding arc of a fixed length $r$.  Its average over
the starting point is $rL$, so it cannot have minimum above $rL$.  Fine
periodic checkerboards make every sliding arc asymptotically have that mass.
Equality means that every sliding arc has exactly its average mass; taking a
distributional derivative says precisely that $\omega$ is invariant under
rotation by $r$.

At times $2\pi m$, all distinct temporal Fourier modes of the circle are
orthogonal.  The full wave Gramian becomes block diagonal by Laplace
eigenspaces, so its bottom eigenvalue is exactly $(T/2)g_1$.  The quantity
$g_1$ is controlled by the even Fourier coefficients of $\chi_\omega$.
Fourier-flat unions of very small cells show that all these coefficients can
be made uniformly small.  Vanishing of every even coefficient, however,
forces $\chi_\omega(x)+\chi_\omega(x+\pi)$ to be constant, which a
$\{0,1,2\}$-valued function permits only when $L=1/2$.

\section*{The geometric functional at every time}

\begin{theorem}\label{thm:g2}
Let $T>0$ and $0<L<1$.  Then
\[
 \sup_{\omega\in\Ucal_L}g_2^T(\omega)=L.
\]
Write $T=2\pi q+r$ with $q\in\N\cup\{0\}$ and $0\leq r<2\pi$.
If $r=0$, every $\omega\in\Ucal_L$ is a maximizer.  If $r>0$, then
$\omega$ is a maximizer if and only if
\begin{equation}\label{eq:invariance}
 \chi_\omega(x+r)=\chi_\omega(x)\quad\text{for almost every }x.
\end{equation}
Consequently maximizers exist for every $L$ when $r/(2\pi)$ is rational, and
no maximizer exists when $r/(2\pi)$ is irrational.  Moreover, for every
$\omega\in\Ucal_L$,
\[
 g_2(\omega):=\lim_{T\to\infty}g_2^T(\omega)=L.
\]
\end{theorem}

\begin{proof}
Put
\[
 F_r(x)=\int_0^r\chi_\omega(x+t)\,dt.
\]
Every complete turn contributes $|\omega|=2\pi L$, and therefore
\begin{equation}\label{eq:turns}
 g_2^T(\omega)=\frac{2\pi qL+\inf_xF_r(x)}{T}.
\end{equation}
Fubini's theorem gives
\[
 \frac1{2\pi}\int_\Torus F_r(x)\,dx=rL.
\]
Thus $\inf_xF_r(x)\leq rL$ (with essential infimum if preferred), proving
$g_2^T(\omega)\leq L$.

For sharpness, divide the circle into $N$ cells of length $2\pi/N$ and retain
the first proportion $L$ of each cell:
\[
 \omega_N=\bigcup_{k=0}^{N-1}
 \left[\frac{2\pi k}{N},\frac{2\pi(k+L)}N\right).
\]
Every arc differs from a union of complete cells only at its two ends.  Hence,
uniformly in $x$ and $0\leq r\leq2\pi$,
\[
 \left|F_r^{\omega_N}(x)-rL\right|\leq\frac{4\pi}{N}.
\]
Together with \eqref{eq:turns}, this proves
$g_2^T(\omega_N)\to L$.

Equality in the minimum--average bound holds exactly when $F_r$ is constant
almost everywhere.  Its distributional derivative is
\[
 F_r'(x)=\chi_\omega(x+r)-\chi_\omega(x),
\]
so equality is equivalent to \eqref{eq:invariance}.  If
$r/(2\pi)=p/s$ in lowest terms, the set
\[
 \bigcup_{k=0}^{s-1}
 \left[\frac{2\pi k}{s},\frac{2\pi(k+L)}s\right)
\]
is invariant under rotation by $r$ and has measure $2\pi L$.  If the ratio is
irrational, ergodicity of irrational circle rotation forces every invariant
measurable set to have measure zero or $2\pi$, impossible here.  Finally,
\eqref{eq:turns} directly gives $g_2^T(\omega)\to L$ as $T\to\infty$.
\end{proof}

\section*{The deterministic constant at resonant times}

\begin{lemma}\label{lem:flat}
For every $0<L<1$, there are sets $\omega_N\in\Ucal_L$ such that
\begin{equation}\label{eq:flat}
 \sup_{\ell\in\mathbb Z\setminus\{0\}}
 \left|\int_{\omega_N}e^{i\ell x}\,dx\right|\longrightarrow0.
\end{equation}
\end{lemma}

\begin{proof}
Partition the circle into $N$ equal cells $I_k$.  Select cells independently
with probability $L$.  The elementary Hoeffding bound, followed by a union
bound over the $N-1$ nonzero discrete Fourier frequencies, shows that one can
choose a realization for which both the discrepancy in the number of selected
cells and every discrete Fourier sum
\[
 \sum_{k\ \mathrm{selected}}e^{2\pi iak/N},\qquad 1\leq a<N,
\]
are $O(\sqrt{N\log N})$.  Add or remove
$O(\sqrt{N\log N})$ cells to leave exactly $\lfloor LN\rfloor$ full cells,
and add a subinterval of one remaining cell to make the measure exactly
$2\pi L$.  The same bounds persist.

For an integer $\ell\neq0$, the integral over the union of full cells factors
as
\[
 \left(\int_0^{2\pi/N}e^{i\ell x}\,dx\right)
 \sum_{k\ \mathrm{selected}}e^{2\pi i\ell k/N}.
\]
The first factor has modulus at most $2\pi/N$; when $N$ divides $\ell$ it is
zero.  The final partial cell contributes at most $2\pi/N$.  Uniformly in
$\ell\neq0$, the result is $O(\sqrt{\log N/N})$, proving
\eqref{eq:flat}.
\end{proof}

\begin{theorem}\label{thm:C}
Let $m\in\N^*$, $T=2\pi m$, and $0<L<1$.  For every
$\omega\in\Ucal_L$,
\begin{equation}\label{eq:exactC}
 C_T(\omega)=\frac{T}{2}g_1(\omega),
\end{equation}
where
\begin{equation}\label{eq:g1fourier}
 g_1(\omega)=L-\frac1{2\pi}\sup_{j\geq1}
 \left|\int_\omega e^{2ijx}\,dx\right|.
\end{equation}
Consequently
\[
 \sup_{\omega\in\Ucal_L}C_T(\omega)=\frac{TL}{2}.
\]
This supremum is attained if and only if $L=1/2$.  For $L=1/2$, the maximizers
are exactly the antipodal selectors
\begin{equation}\label{eq:selector}
 \chi_\omega(x)+\chi_\omega(x+\pi)=1\quad\text{a.e.}
\end{equation}
\end{theorem}

\begin{proof}
The positive circle frequencies are the integers $j\geq1$, with eigenspaces
\[
 E_j=\operatorname{span}\{\cos(jx),\sin(jx)\}.
\]
In the source's
$y^+,y^-$ variables, integration over $[0,2\pi m]$ kills every temporal term
with nonzero integer frequency.  The observation Gramian is therefore block
diagonal: on each of the two time branches and each $E_j$, its quadratic form
is $T\int_\omega|u|^2$, while the source normalization contributes the factor
$1/2$.  Taking the lowest block eigenvalue proves \eqref{eq:exactC}.

Every unit vector in $E_j$ is a phase shift of
$\pi^{-1/2}\sin(jx)$.  Since
\[
 \sin^2(jx+\theta)=\frac12-\frac12\cos(2jx+2\theta),
\]
minimization over $\theta$ gives \eqref{eq:g1fourier}.  Lemma~\ref{lem:flat}
then yields $\sup g_1=L$ and hence the stated supremum of $C_T$.

If $g_1(\omega)=L$, every even nonzero Fourier coefficient of
$\chi_\omega$ vanishes.  On $[0,\pi)$ define
\[
 h(x)=\chi_\omega(x)+\chi_\omega(x+\pi).
\]
All nonconstant Fourier coefficients of $h$ vanish, so Fourier uniqueness
gives $h=2L$ almost everywhere.  But $h$ takes values only in
$\{0,1,2\}$; as $0<L<1$, this forces $L=1/2$ and $h=1$.  Conversely,
\eqref{eq:selector} makes every even nonzero Fourier coefficient vanish and
therefore gives $g_1=L$.  This proves the classification.
\end{proof}

\section*{Upgrade attempts, limitations, and novelty bounds}

The finite-time $g_2^T$ theorem is a complete result for the circle and
strictly strengthens its long-time $g_2$ subcase.  The deterministic theorem
settles an infinite family of observation times but not arbitrary $T$.

An additional upgrade attempt used rotation averaging: concavity and
translation invariance give the natural relaxed upper bound obtained from the
constant density $L$.  The checkerboards in Theorem~\ref{thm:g2} also make the
high-frequency geometric obstruction converge uniformly to this relaxed
value.  What remains for arbitrary $T$ is a two-scale estimate controlling the
off-diagonal wave Gramian when the observation cells and the Fourier modes
both tend to high frequency.  Weak-* convergence of characteristic functions
does not control this operator norm, so the resonant proof cannot simply be
perturbed.  Higher-dimensional manifolds likewise lose the one-dimensional
sliding-arc reduction.

On 11 August 2026, searches covered the exact source title and question,
$g_2^T$/circle/sliding-window formulations, resonant-time deterministic
observability, and the citing literature.  Privat--Tr\'elat--Zuazua
\cite{ptz} already note Parseval equality of deterministic and randomized
constants at common periods in a one-dimensional interval and establish
related no-gap/nonattainment results for the randomized spectral design.
No source located in the bounded search states Theorem~\ref{thm:g2} or this
combined circle classification.  Novelty confidence is moderate: the proofs
are elementary once the circle reduction is made, and the resonant part has
clear antecedents.

The included script numerically checks sliding-window convergence,
Fourier-flat finite-cell samples, and the formulas for $g_1$; it is not used in
the proof.  Human review should focus on the normalization in
\eqref{eq:exactC}, the uniform-in-frequency step in Lemma~\ref{lem:flat}, and
the distinction between essential and pointwise infima for arbitrary
measurable representatives.

\begin{thebibliography}{9}
\bibitem{source}
E. Humbert, Y. Privat, and E. Tr\'elat,
\emph{Observability properties of the homogeneous wave equation on a closed
manifold}, arXiv:1607.01535, concluding shape-optimization problems (PDF p.~22).

\bibitem{ptz}
Y. Privat, E. Tr\'elat, and E. Zuazua,
\emph{Optimal observability of the multi-dimensional wave and Schr\"odinger
equations in quantum ergodic domains}, arXiv:1211.5655.

\bibitem{hoeffding}
W. Hoeffding,
\emph{Probability inequalities for sums of bounded random variables},
J. Amer. Statist. Assoc. 58 (1963), 13--30.
\end{thebibliography}

\end{document}
